% ============================================================
%  Chapter 5 — Conclusion
% ============================================================
\chapter{Conclusion}
\label{ch:osszegzes}

\section{Main Findings}

KCMamba integrates a differentiable Kalman filtering mechanism into the Mamba
SSM. The experiments on three datasets lead to the following conclusions:

\begin{enumerate}
    \item \textbf{Noise robustness.} Going from $\sigma=0$ to $\sigma=5$,
          KCMamba's average MSE degradation is $+1\,528\%$, versus
          $+28\,234\%$ for Fair Mamba, an $\mathbf{18.5\times}$ difference.
          In the Exchange/$s=16$ setting this becomes five orders of magnitude:
          $+66\%$ vs.\ $+152\,350\%$.

    \item \textbf{Kalman gain adaptation.} The model learned on its own to
          reduce $K$ as noise increases ($0.706 \to 0.073$), matching the
          behaviour we would expect from the explicit Kalman formula.
          The architectural prior genuinely guides learning.

    \item \textbf{Scarce data.} With stride $s=16$, the Kalman structure acts
          as an implicit regulariser that partially compensates for the lack of
          training windows.
\end{enumerate}

\section{Limitations}

KCMamba does not win in every configuration:
\begin{itemize}
    \item At $\sigma=0$, $s=16$, Fair Mamba nearly memorises and achieves
          much better MSE.
    \item On ETTh1/$s=1$, LSTM is the best; its gating mechanism exploits
          the large data volume effectively.
    \item Only 5 training epochs were run, which is not necessarily a converged
          state. Longer training might reorder the results.
    \item Only a seq=96/pred=96 horizon was evaluated; other horizons may show
          different rankings.
\end{itemize}

\section{Future Work}

\begin{itemize}
    \item A systematic sweep over $s \in \{1, 4, 8, 16, 32, 64\}$ would
          identify the data-scarcity threshold at which the Kalman
          structure's regularising effect first emerges.

    \item At prediction horizons of 336, 720, or 960 steps, Mamba's
          $O(T \log T)$ complexity may confer a meaningful advantage
          over $O(T^2)$ Transformer baselines.

    \item Replacing simulated Gaussian perturbations with sensor
          dropouts, missing values, and outliers would bring the
          evaluation closer to real-world deployment conditions.

    \item Introducing a meta-network that estimates $\sigma$ as a
          trainable parameter would allow the model to adapt
          automatically to environments with unknown noise levels.

    \item Replacing the scalar $Q_t$, $R_t$ with a diagonal Cholesky
          factorisation would yield richer uncertainty estimates and
          potentially improve calibration.
\end{itemize}

\section{Closing Remarks}

KCMamba demonstrates that classical Kalman filtering principles and modern
sequence models are not mutually exclusive. The explicit noise model provides
a strong inductive bias, particularly where data is scarce or noise is high.
The trading system described in Chapter~\ref{ch:implementacio}
shows that this model is not just an experimental prototype but a solution
integrable into a real trading environment.
